\clearpage
\section*{Abstract}
\addcontentsline{toc}{section}{Abstract}

This study developed a residential property valuation framework for Metro Cebu using machine learning and GIS-derived spatial features. The work addressed the gap between administrative benchmarks, lender-driven appraisals, and market-facing price signals by assembling a hybrid dataset of residential properties from institutional acquired assets, foreclosure listings, and open-market online listings. The study focused on six local government units in Metro Cebu: Cebu City, Mandaue City, Lapu-Lapu City, Talisay City, Minglanilla, and Consolacion.

The methodology combined structural, locational, administrative, and geospatial variables. Property records were geocoded, matched with barangay-level BIR zonal values, and enriched with accessibility and neighborhood-context features, including proximity to major economic nodes, amenity access, and a neighborhood price lag. Because property types are priced on very different per-square-meter logics, separate models were fitted for condominiums, houses, and vacant lots, on a price-per-square-meter target. Within each stratum, three predictive approaches were compared: hedonic regression, Random Forest, and XGBoost. Under leak-free grouped cross-validation, the tree-based models outperformed the hedonic baseline, while Random Forest and XGBoost performed comparably; Random Forest was deployed in all three strata for its robustness on small samples and its simpler implementation. The deployed models reached a typical error of about 19, 23, and 38 percent (median absolute percentage error) for condominiums, houses, and vacant lots, with location features carrying most of the predictive signal. Model outputs were interpreted through SHAP-based explainability and prepared for use through a QGIS map and a web-based decision-support application centered on a predicted residential price surface.

The manuscript positions the study as a predictive and prescriptive tool for residential valuation practice in Metro Cebu. Its contribution lies in combining property-level machine learning, spatial feature engineering, and transparent explanation into a locally grounded decision-support framework that is more responsive to urban change than static administrative benchmarks alone.
